% 第 33 章　把光当成射线
\chapter{把光当成射线}

\step{反常识开场}

请你现在做一件事：找一根筷子，插进半杯清水里，从侧面看它。筷子在水面处“断”了——水下那一截像被人掰了一下，错位地折向一边。把筷子拔出来，它完好无损；再插回去，它又“断”了。

你当然不会真的以为水有什么神力把筷子掰弯了。可问题来了：你的眼睛一辈子信奉“看到的东西就在视线的方向上”，这一次它却指着错误的方向信誓旦旦。水下那截筷子发出的光，明明到达了你的眼睛，你的眼睛却把它“安放”在了一个错误的位置。光在从水到空气的路上，一定发生了某件事。

这个“骗局”不只是餐桌上的把戏，它影响过真本事。生活在江边海边的人用鱼叉叉鱼，代代相传的口诀是：不要瞄准你看见的那条鱼，要瞄准它的下方。孩子学叉鱼，前几叉总是落空——因为他瞄准的是眼睛报告的位置，而鱼根本不在那里。有经验的渔人不谈物理，却早已向水面的骗局缴了学费。

更离奇的事情发生在水面之下。如果你潜到泳池底，仰面往上看，只要你的位置足够深、看的方向足够斜，你会发现头顶的水面不是透明的窗，而是一面镜子——岸上的灯、天花板的倒影清清楚楚，水面之外的景物却完全消失。同一片水面，从上看是窗，从下看是镜。窗户什么时候会变成镜子？判据是什么？

再看一个廉价却耐玩的现象。夏天的柏油马路被晒得滚烫，开车时远看前方，路面上像有一汪汪积水，倒映着天空和对向的车灯；开近了，水洼消失，远处又出现新的水洼。司机知道那是假的，可它看起来那么像水的反光。沙漠旅人看到的“湖泊”，也是同一个东西。根本没有水面，光却像是被一面不存在的镜子反射了回来。

\begin{mainline}
这一章的全部麻烦可以压成一个问题：光在路上是走直线的（影子、针孔相机都是证据），可它一旦越过两种物质的边界，就会突然拐弯，拐多少还有精确的规律。我们要弄清：这条“拐弯的规则”是什么，它又服从什么更深的道理。
\end{mainline}

\step{先猜一猜}

在翻下页之前，请先写下你自己的猜测：筷子在水面处“折断”，最可能的原因是什么？

第一种猜测：是眼睛（或大脑）的错觉。水本身对光没有影响，只是人脑不习惯透过水看东西，自作主张地补出了一个错误的画面。

第二种猜测：水像一种“阻力”，把光拖慢了，光走得费劲，所以歪了。就像人涉水走路会踉跄一样。

第三种猜测：光在两种物质的分界面上，会按照某种确定的规则改变方向。偏转不是错觉，也不是含糊的“阻力”，而是可以测量的、每次都一样的角度变化。

也请你顺便预测一件后面要检查的事：如果让光从空气斜射进水里，你猜它是向竖直方向靠拢，还是远离竖直方向？反过来，从水里射向空气呢？把你的答案记下来。

还有一个更细的猜测题。斜着射入水面的光，是一部分进入水中、一部分被水面反射，还是只能发生其中一种？如果你拿不准，就去湖边看一次：黄昏时看湖面，你既能看见天空的倒影，又能看见水下的石头——两样同时在。这个事实待会儿会帮我们排除掉一些模型。

\step{动手实验}

\begin{experiment}[消失的硬币]
找一个不透明的碗或马克杯，把一枚硬币放在碗底。你站好位置，慢慢后退（或者把碗慢慢往前挪），直到碗沿刚好挡住硬币——硬币刚刚看不见为止。保持头和碗都不动。

现在请另一个人缓缓往碗里倒水，你保持原来的姿势盯着碗底看。随着水面升高，刚才被挡住的硬币会一点一点“浮”回你的视野，最后整枚都能看见。倒掉的如果不是水而是别的透明液体（比如食用油），“浮回来”的程度还不一样。
\end{experiment}

这个实验里，硬币没有动，你的眼睛没有动，碗也没有动，唯一变化的是碗里多了水。结论逃不掉：水把从硬币射向你眼睛的光，在水面处掰了一个弯。而且硬币是从碗沿方向“浮”出来的——这说明光折弯有确定的方向性，不是乱折。

如果你家有激光笔（玩的时候记住：永远不要直射眼睛），可以再做一组观察。在透明方盒里装上水，滴入一两滴牛奶搅匀，让光路显形。从空气斜着把激光射向水面，你会同时看到三件事：一条光从水面弹回空气（反射）；一条光钻进水里，但在水面处明显折了一下（折射）；折射那条比入射那条更靠近竖直方向。把水换成糖水、油，折的程度会变化。

请回答三个观察题：入射光越斜，折射光偏离原方向越多还是越少？垂直水面直射时，光还拐弯吗？把盒子举高，让激光从水里向上斜射向水面，不断增大倾斜程度，会发生什么？第三问先记住现象，本章后半段它是主角。

\step{旧模型遇到麻烦}

先替旧模型说句公道话。“光沿直线传播”是一条非常成功的经验。它解释手影：手挡住光，身后光到不了的地方就是暗的；它解释日食：月球正好挡在太阳和地球之间，本影里的人看到太阳被整个遮住；它解释针孔相机：物体上每一点发出的光穿过小孔，在后面的屏上落在一个确定的位置，于是投出倒立的像。工匠用它校准直线，测量员用它瞄准，古人用它立竿测日影定节气。在这套语言里，光是一种沿直线飞的东西，整个光学就是画直线。

可是，把这套语言拿到水面，它立刻撞墙。

第一处撞墙：折射。“光走直线”没有为“分界面”预留任何角色。筷子实验和硬币实验里的偏转是确定、可重复的，不是模糊的边缘效应。旧模型只能说“光拐弯了”，说不出拐多少、朝哪拐、为什么玻璃拐得比水多。

第二处撞墙：它连自己最擅长的反射也说不清。人人都知道“入射角等于反射角”，可是为什么相等？“走直线”只是描述了路线，没有给出任何理由让两个角非得相等不可。就像你记录了一千次路口的行车轨迹，却不明白司机为什么都那样转弯——有数据，没有道理。

第三处撞墙，也是最微妙的：旧模型说不清楚“光走直线”到底是实验事实，还是模型的假设。严格讲，它是有条件的近似：光在同一种均匀物质里确实走直线，可一旦物质变了（空气变水）、或者物质本身不均匀（被晒热的马路上空的空气），直线就保不住了。把一条有条件的近似当成无条件的定律，正是直觉在后面要付的学费。

\begin{mainline}
旧模型的账目可以总结为：它准确\textbf{描述}了三件事——直线传播、反射、折射，却一件也\textbf{解释}不了。我们需要一个新的概念工具，既能装下这三条描述，又能回答“为什么偏偏是这个角度”。
\end{mainline}

\step{发明新概念}

我们先做一个思想上极为重要、又极易被误解的抽象：\textbf{光线}。

在纸上画一条带箭头的细线，代表光前进的方向和路线——这条线叫光线。它\textbf{像什么}：像地图上给登山者画的路线路径，告诉你从哪里出发、朝哪个方向、在哪里转折。它\textbf{不像什么}：它不是一根真实存在的细丝。你用激光笔看到的那条亮线，是光路上灰尘和水汽把一小部分光散射进你眼睛的结果；真空中前进的光，从侧面看是完全看不见的。光线也不是“光的实体”，你不能把它剪断、称重量。它只是我们给光的前进方向起的一个几何名字。这个抽象之所以强大，是因为影子、反射、折射全部可以翻译成关于“线”的几何问题，而几何问题是可以画图、量角、计算的。

接下来把三个现象翻译成几何语言。在分界面上作一条垂直于界面的参考线，叫\textbf{法线}。入射光线与法线的夹角叫\textbf{入射角}，反射光线与法线的夹角叫\textbf{反射角}，进入另一种物质的光线与法线的夹角叫\textbf{折射角}。三个角都从法线量起，不是从界面量起——这个约定是为了让定律写出来最简洁。

\begin{center}
\begin{tikzpicture}[scale=1.0]
  % 界面与法线
  \draw[thick] (-4,0) -- (4,0);
  \draw[dashed] (0,-2.6) -- (0,2.6) node[above] {法线};
  \node at (-3.6,0.35) {空气};
  \node at (-3.6,-0.35) {水};
  % 入射光线
  \draw[->,thick] (-3,2.4) -- (0,0) node[midway,left=2pt] {入射光};
  % 反射光线
  \draw[->,thick] (0,0) -- (3,2.4) node[midway,right=2pt] {反射光};
  % 折射光线（折向法线）
  \draw[->,thick] (0,0) -- (1.5,-2.4) node[midway,right=2pt] {折射光};
  % 角标注
  \draw (0,1.1) arc (90:128.7:1.1) node[pos=0.5,above=3pt] {$\theta_1$};
  \draw (0,1.7) arc (90:51.3:1.7) node[pos=0.5,above=3pt] {$\theta_1'$};
  \draw (0,-1.3) arc (270:302:1.3) node[pos=0.55,below=4pt] {$\theta_2$};
\end{tikzpicture}
\end{center}

然后是两条从大量测量里归纳出来的实验事实：

\textbf{反射定律}：反射光线在入射光线和法线决定的平面内，且反射角等于入射角，\(\theta_1'=\theta_1\)。这条定律对镜子、水面、地板、任何光滑界面都一样，精确得可以拿去做卫星测距。日常里它的出场率远超想象：潜望镜靠两面平行的镜子把光折两次，让战壕和水下的人越过障碍观察；公路两侧的反光标志和自行车尾灯里的“猫眼”，内部是许多微小的直角反射结构，能把车灯的光几乎原路送回司机眼里——逆着来路反射回去，所以夜里格外刺眼。

\textbf{折射定律}：折射光线也在同一平面内，但折射角一般不等于入射角。同一个入射角下，从空气进水的偏折程度是固定的，进玻璃的偏折程度是另一个固定值。也就是说，每种透明物质带着一个属于自己的“弯折能力”的数字。我们给这个数字起名为\textbf{折射率}，记作 \(n\)。真空（近似地说，空气）的折射率规定为 \(1\)；水约为 \(1.33\)，普通玻璃约 \(1.5\)，钻石高达 \(2.42\)。折射率越大，界面把光掰弯得越狠。

折射率还有一个更深刻的物理身份。实验上可以直接测量光在不同物质中的传播速度：真空中光速 \(c\approx 3.0\times10^8\) 米每秒，而在水中只有约 \(2.25\times10^8\) 米每秒，在玻璃里约 \(2.0\times10^8\) 米每秒。测量发现，对每一种透明物质，
\[
v=\frac{c}{n},\qquad\text{即}\qquad n=\frac{c}{v}.
\]
折射率就是“真空光速与介质中光速之比”。这是一个值得停下来看的等式：一个描述“光拐多大弯”的几何数字，竟然同时就是“光在这种物质里走多慢”的数字。拐弯和变慢，是同一件事的两个侧面。这一章我们先接受这个实验事实；为什么物质会同时决定这两件事，要等第 35、36 章把光当成波、看清光与物质的相互作用之后，才能真正回答。

到这里，反射和折射仍然是两条独立的经验规则。它们有没有共同的根？十七世纪的费马提出了一个大胆得近乎离奇的想法：\textbf{光从一点走到另一点，走的总是所花时间取稳定值（通常是最短）的那条路。}

先别急着觉得玄。想象海滩上的救生员：溺水者在海里的斜前方，救生员在沙滩上跑得很快、在水里游得很慢。他该沿直线冲过去吗？不该。聪明的做法是在沙滩上多跑一段、瞄准一个偏上游的入水点，让“慢速的水中路”尽量短。最快的路线不是直线，而是在岸边“折”了一下的折线——而且两种速度差别越大，折得越明显。这正是光在水面上的行为：空气里的“沙滩”上光速快，水里“游得慢”，于是光选择了一条在界面处折向法线的路线。

这个比喻\textbf{像的部分}：两种速度、两种区域、在边界上择优，结构与折射完全同构。它\textbf{不像的部分}：救生员有意识、会比较、会决策；光没有。“光选择最省时的路”不是说光在出发前把所有路线试了一遍再挑——这是数学模型的说话方式，不是光的内心戏。至于为什么一条没有时间观念的物理定律会以“时间最短”的面目出现，要到第 35 章的波动图景里才会露出端倪：那时我们会看到，费马原理是波在波长极短时的自然结果。

费马原理还顺手送了一件礼物：\textbf{光路可逆}。既然“时间最短”是从 \(A\) 到 \(B\) 这条路径本身的性质，那么把光从 \(B\) 沿同一条路射回去，这条路依然是时间最短的路。所以你若能在水面上方看见潜水员的眼睛，潜水员也一定能看见你——互相看见这件事，从来不需要两个人各自“发射视线”，只有真实的光在路上跑。光路可逆看起来是句废话，却是后面设计复杂光学仪器时省掉一半计算的关键。

顺便把影子也重新画一遍，因为它马上要在边界一节当证人。点光源（比如激光笔照一个小孔）照出的影子边缘锐利；而太阳、日光灯管这类有尺寸的光源照出的影子分两层：所有方向的光都被挡住的区域叫本影，只有一部分光被挡住的区域叫半影。日食就是月球把这套影子投在地球上：站在本影里的人看到日全食，站在半影里的人看到日偏食。这一套完全用“直线传播”就能画出来——它是射线模型最漂亮的胜仗，记住它成立的前提：光源、障碍物的尺寸都远大于光的波长。

\step{一句话模型}

\begin{oneline}
把光当成沿“光线”前进的模型：在均匀物质里光走直线，在两种物质的界面上，一部分光按“反射角等于入射角”弹回，另一部分按斯涅尔定律折入；而这两条规则背后是同一条更深的账——光走花时间取稳定值的路（费马原理），折射率 \(n=c/v\) 记录每种物质让光走多慢。
\end{oneline}

这句话里有三个层次，请分清：哪些是实验事实（直线传播、反射定律、折射定律、介质中的光速），哪些是数学模型（光线、角度、折射率、光程），哪些是解释框架（费马原理）。把事实、模型、解释混在一起，是光学里绝大多数糊涂观念的来源。

\step{数学翻译器}

现在把直觉翻译成可以算的式子。

\textbf{反射}的式子简单到几乎不像公式：
\[
\theta_1'=\theta_1 .
\]
用特殊情况检查它。垂直照射：\(\theta_1=0\)，则 \(\theta_1'=0\)，光原路返回——拿手电垂直照镜子，光斑回到手边，对。掠射：\(\theta_1\to 90^\circ\)，反射光也贴着界面飞出——打水漂的石片在极低角度下会被水面“完美”弹起，日落时低角度阳光在水面上反射极强，对。方向反转：把光沿原反射光路逆着射回去，它沿原入射光路出来——光路可逆，对。

\textbf{折射}的式子是斯涅尔定律：
\[
n_1\sin\theta_1=n_2\sin\theta_2 .
\]
读一下式子的两边：左边是“入射侧”的物质和角度，右边是“折射侧”的；它说的是：\(n\sin\theta\) 这个组合，在穿越界面时保持不变。越“慢”的物质（\(n\) 大），角就得越小来补偿——光在慢物质里更贴近法线。

逐项用特殊情况检查。

检查一：垂直入射。\(\theta_1=0\)，左边为零，所以 \(\theta_2=0\)：垂直射入不偏折。这和动手实验里激光直射水面不拐弯一致。但请注意一个容易滑过去的陷阱：方向没变，不等于什么都没变——光速从 \(c\) 变成了 \(c/n\)，穿过同样厚度花的时间变长了。眼睛只对方向敏感，所以你看不见这次变化，但光自己已经“记账”了。

检查二：两种物质相同。\(n_1=n_2\)，则 \(\theta_2=\theta_1\)，不偏折——本来就没有界面，当然不折。一个漂亮的应用：把玻璃棒插进折射率与玻璃几乎相同的油里，玻璃棒会“隐身”——界面上没有弯折，光线直接穿过，眼睛找不到玻璃在哪儿。

检查三：从空气进水。\(n_1=1\)，\(n_2=1.33\)，则 \(\sin\theta_2=\sin\theta_1/1.33\)。取 \(\theta_1=45^\circ\)：\(\sin\theta_2\approx 0.707/1.33\approx 0.53\)，\(\theta_2\approx 32^\circ\)。折射光明显更靠法线，与实验中看到的弯折方向一致——你之前猜对了吗？

检查四：从水出空气，并把入射角不断调大。这时 \(n_1=1.33\)，\(n_2=1\)，式子要求 \(\sin\theta_2=1.33\sin\theta_1\)。可是正弦最大只能到 \(1\)：当 \(\theta_1\) 大到使 \(1.33\sin\theta_1=1\) 时，折射角已经达到 \(90^\circ\)，再增大入射角，方程\textbf{无解}。数学在报警：折射光消失了。发生报警的临界角度满足
\[
\sin\theta_c=\frac{n_2}{n_1}\quad(n_1>n_2).
\]
水对空气：\(\sin\theta_c=1/1.33\)，\(\theta_c\approx 48.6^\circ\)。超过这个角，全部光都被界面反射回来，一丝不漏——这就是\textbf{全反射}。这里就是一个会骗到直觉的反例：多数人的直觉是“界面总会透过去一点，只是多少而已”，毕竟玻璃既能反光又能透光。但全反射说不：一旦越过临界角，透射不是“变少”，而是干脆为零，反射率跳到完美的百分之百——比任何镜子都彻底（镜子还要吸收百分之几）。

全反射是光纤的命根子。一根比头发粗不了多少的玻璃丝，芯的折射率略高于外皮。光在芯里以足够斜的角度前进，每撞一次壁就被全反射弹回，像在镜子的隧道里“之”字形前进，几公里、几十公里都不漏出去。海底光缆就是这样把整部电影的信号从一个大洲送到另一个大洲；医生用的内窥镜、胃镜，也是靠一束光纤把光送进人体、再把图像传出来。

\begin{center}
\begin{tikzpicture}[scale=1.0]
  % 光纤
  \draw[thick] (0,0.8) -- (8,0.8);
  \draw[thick] (0,-0.8) -- (8,-0.8);
  \node at (0.2,1.15) {外皮（$n$ 较小）};
  \node at (0.6,0) {芯（$n$ 较大）};
  % 之字形光线
  \draw[->,thick] (-0.8,-0.4) -- (0,0) ;
  \draw[thick] (0,0) -- (1.2,0.8) -- (3.0,-0.8) -- (4.8,0.8) -- (6.6,-0.8) -- (7.8,0);
  \draw[->,thick] (7.8,0) -- (8.6,0.3);
  % 反射点标记
  \foreach \x/\y in {1.2/0.8, 3.0/-0.8, 4.8/0.8, 6.6/-0.8}
    \fill (\x,\y) circle (0.05);
  \node at (4.8,1.2) {每次碰壁：入射角大于临界角，全反射};
\end{tikzpicture}
\end{center}

\textbf{看东西的深浅}也能算了。从水面上方竖直向下看，水下深度为 \(h\) 的物体，看起来的深度（视深）是
\[
h'=\frac{h}{n}.
\]
用真实数据估一笔账：泳池水深 \(2.0\) 米，\(n=1.33\)，则 \(h'\approx 2.0/1.33\approx 1.5\) 米。也就是说，眼睛报告的深度比真实深度浅了四分之一——这就是为什么孩子看见“才到腰深”的水，跳下去会没顶；也是叉鱼口诀“瞄准鱼的下方”的定量版。这个估算用的是近似条件（视线接近竖直），推导见下面的数学桥盒子。

\begin{mathbridge}
\textbf{从费马原理推出斯涅尔定律。}设光从空气中一点 \(A\) 出发，经界面上一点 \(P\) 到水中一点 \(B\)。把 \(A\)、\(B\) 到界面的垂足距离分别记为 \(a\)、\(b\)，垂足间距固定为 \(L\)，设 \(P\) 距 \(A\) 的垂足为 \(x\)。总时间
\[
T(x)=\frac{n_1}{c}\sqrt{a^2+x^2}+\frac{n_2}{c}\sqrt{b^2+(L-x)^2}\,.
\]
费马原理说真实的 \(x\) 使 \(T\) 取稳定值，即 \(\mathrm{d}T/\mathrm{d}x=0\)：
\[
n_1\frac{x}{\sqrt{a^2+x^2}}=n_2\frac{L-x}{\sqrt{b^2+(L-x)^2}}\,.
\]
左边恰好是 \(n_1\sin\theta_1\)，右边恰好是 \(n_2\sin\theta_2\)。斯涅尔定律从“时间取稳定值”一个要求里完整掉了出来，连符号都对。

\textbf{视深公式的推导。}从水下物体发出两条靠得很近的光，一条竖直向上，一条与法线成小角 \(\theta_2\)，出水后按斯涅尔定律变成 \(\theta_1\approx n\,\theta_2\)（小角近似）。眼睛把出射光沿直线延长回去，交点在深度 \(h'=h\tan\theta_2/\tan\theta_1\approx h/n\) 处。注意这里用了小角近似，所以斜着看时视深公式要修正——这正是侧看泳池底变形的原因。
\end{mathbridge}

\begin{challenge}
\textbf{“最短”其实应该是“稳定”。}费马原理的准确版本是：光程（传播时间乘以 \(c\)，即 \(\int n\,\mathrm{d}s\)）取\textbf{稳定值}——可以是最小，也可以是最大或马鞍形的拐点。一个干净的反例是椭球面镜：从焦点 \(F_1\) 发出的光经镜面反射到焦点 \(F_2\)，所有路径的光程严格相等（椭圆的定义就是距离之和为常数），每条都是稳定值，没有谁是“最短”。更一般的界面形状下，稳定路径还可以是光程极大。“最短时间原理”这个流行叫法在多数日常场景碰巧成立，但作为定律是错的——稳定，才是费马真正说的事。

\textbf{和力学的隐秘通道。}如果你学过第 13、14 章，会发现这里的味道似曾相识：力学里质点走使作用量取稳定值的轨道，光学里光走使光程取稳定值的路径。这不是巧合。哈密顿最早看出，力学和几何光学在数学上是同一个结构——粒子的轨道之于力学，正如光线之于光学。这条通道后来成为薛定谔猜出波动力学的向导之一，第 43 章还会走到它的另一头。
\end{challenge}

\step{模型的边界}

射线模型是一张极好用的地图，但地图上没画的东西，走真路时会绊倒你。

\textbf{边界一：障碍物不能太小。}光线模型默认光严格沿直线走，影子边缘绝对清晰。可如果你盯着日光灯下铅笔的影子仔细看，影子边缘并不是刀切的，而是有一层由明到暗的过渡；让光穿过足够细的缝，它不但不沿直线前进，反而会散开。这些现象在障碍物的尺寸小到和光的波长可比拟（大约几百纳米，比头发丝的百分之一还细）时必然出现。射线模型对此完全失语——这不是模型的错误答案，而是模型根本没被问过这个问题。为什么光会绕弯、缝越窄散得越开，是第 35 章波动光学的主菜。那里我们会看到：射线不是错的，而是波长趋于零时的极限图景。

\textbf{边界二：光线不是实体，也不解释“多少”。}本章的定律只管方向，不管能量分配。为什么水面斜看时倒影特别亮、正看时几乎透明？多少光反射、多少光折射，随角度和偏振怎么变？这需要第 36 章的菲涅耳公式。只用射线模型讨论“光强”，是在地图上量海拔。

\textbf{边界三：折射率依赖于颜色。}我们说“水的折射率是 \(1.33\)”，其实那是黄绿光的数值；紫光在水中的折射率比红光略大。所以白光穿过棱镜会被拆成彩色光带，所以钻石的火彩五颜六色。射线模型本身能容忍“每种颜色一个 \(n\)”，但它解释不了为什么 \(n\) 会变——那需要知道光和物质内部电荷的相互作用，同样留到第 35、36 章。

\textbf{边界四：常见的误用。}一是把“光走直线”当成无条件的宇宙法则——它只在均匀介质里成立，海市蜃楼就是反例（下面细说）。二是把介质中光速变慢想象成“光子被原子撞来撞去、绕了远路”——这不是正确的微观图景；真实图像是入射电磁场驱动物质中的电荷振动，电荷再辐射，叠加出来的合成波以 \(c/n\) 的相速度前进，细节在第 36 章。三是把费马原理当成“光会思考”的证据——它是极有效的计算规则，其深层理由在波动图景中，不是目的论。

最后交代一下射线模型和已知旧知识的衔接。第 25 章说过，光是一种电磁波。那一章描述的是电场磁场在空间中的传播，这一章画的却是细线，两种说法怎么并存？答案是：射线模型就是电磁波在“波长比所有障碍物都短得多”时的速写。当波长短到一切东西在它眼里都是庞然大物时，波的行为收敛成几何的线——正如大比例尺地图上，蜿蜒的河流收敛成一条曲线。所以本章不是推翻了“光是电磁波”，而是找到了它最省事的使用方式；它的失效边界，也正是波动细节开始显形的地方。

\step{回头解释开场现象}

现在盘点开场时欠下的四笔账。

\textbf{筷子为什么“折断”。}水下筷子上每一点发出的光，出水时按斯涅尔定律折离法线（从水到空气，\(n\) 由大变小，角由小变大）。你的眼睛和大脑只会做一件事：把收到的光沿直线往回延长。延长线汇聚的位置，比筷子的真实位置更浅、更靠近你——于是水下部分整体“上浮”且错位，水面处出现折角。这不是眼睛的错觉，而是眼睛忠实地执行了“光走直线”的规则，恰好这条规则在水面之上是对的、之下是错的。眼睛没有被骗，它只是不知道光路上有一段是在水里走的。

\textbf{水面何时从窗变成镜子。}潜水员向上看，视线与水面的夹角决定了入射角。当入射角小于临界角 \(48.6^\circ\)，水面是窗，能看到岸上的世界；大于临界角，全反射启动，水面变成完美的镜子，映出水下的景物。所以潜水员头顶有一个圆形的“亮窗”（斯涅尔窗），窗外整个岸上的半球世界被压缩在这个圆里，圆之外全是镜中的水下倒影。

\textbf{马路上的“水洼”。}烈日下，贴地几厘米的空气被路面烤热，热空气稀薄，折射率比上方冷空气略低——空气变成了一种折射率从下往上连续增大的渐变介质。来自天空的光向下斜射时，每穿过一层更“快”的空气，就被掰得离法线更远一点，路径连续弯曲，最终弯到超过局部的临界条件，掉头向上，进入司机的眼睛。司机把光沿直线延长回去，在路面上“看到”了天空的倒影——大脑坚称那里有水面，因为它只认得“镜面反射”这一种能让天空出现在路面上的方式。没有水面，是一整层渐变空气合伙扮演了镜子。这就是海市蜃楼的下现蜃景；海面或极地冷空气层上的上现蜃景（景物悬在空中）原理相同，只是温度梯度反过来。顺便说，海市蜃楼也是“光在均匀介质中走直线”这条近似被自然界公开违反的日常案例——它每天都发生，只是多数人把证据当成了幻觉。

\textbf{被大气“抬”出来的太阳。}渐变折射率的把戏每天都在地平线上演，只是幅度小得多。整层大气的折射率从太空的 \(1\) 连续增大到海平面的约 \(1.0003\)，来自地平线方向的星光、阳光在向下穿过大气时被一点点掰弯。结果：当你看到太阳的下边缘刚好贴上地平线时，太阳的真实位置其实已经完全在地平线之下——它被大气折射“抬”高了约半度，恰好约等于一个太阳圆面的直径。这笔账换算成时间很实在：地球转过半度需要约两分钟，也就是说大气每天替我们“偷”来早晨和傍晚各约两分钟的白昼。天文学家表格里的日出日落时间，都必须先算上这一笔折射修正才和钟面对得上。这是渐变折射率最安静、也最普遍的一次表演：不需要沙漠，不需要热浪，只要有一颗行星裹着大气。

\begin{center}
\begin{tikzpicture}[scale=1.0]
  % 渐变空气层
  \foreach \y in {0.5,1.0,1.5}
    \draw[thin,gray] (0,\y) -- (9,\y);
  \draw[thick] (0,0) -- (9,0) node[midway,below=2pt] {滚烫的路面};
  \node[gray] at (9.6,1.55) {冷空气 $n$ 大};
  \node[gray] at (9.6,0.3) {热空气 $n$ 小};
  % 弯曲光线
  \draw[thick,->] (0,3.0) .. controls (3,0.9) and (4.5,0.55) .. (5.5,0.75);
  \draw[thick,->] (5.5,0.75) .. controls (6.8,1.05) and (8,1.9) .. (9,2.7);
  \fill (0,3.0) circle (0.06) node[left] {天空来的光};
  \fill (9,2.7) circle (0.06) node[right] {司机的眼睛};
  % 虚像
  \draw[dashed] (9,2.7) -- (4.3,0) ;
  \node at (4.3,-0.7) {大脑把光沿直线延长：以为路面有水面};
\end{tikzpicture}
\end{center}

\textbf{激光笔的光柱。}那个顺带的问题：为什么能看见光柱？因为空气里的尘埃把一小部分光散射向各个方向，其中一些进了你的眼睛。在真正干净的真空里，一束光从面前经过，你什么也看不到。看见光柱这件事，恰恰是光“没有”完全沿直线走进你眼睛的证据。

\step{脑洞实验与挑战题}

\begin{thought}[让光绕地球一圈]
光沿直线走，地球是圆的，所以光贴着地平线射出去，会一去不复返。可是渐变折射率给了我们一个疯狂的选项：如果存在一层包裹整个星球的大气，折射率随高度平滑下降，下降得恰到好处，能不能让一束光的路径连续下弯，弯成的圆弧正好和地球表面平行——光绕地球一整圈，回到你的后脑勺？

理论上，只要折射率梯度精确满足条件（粗略地说，光路曲率等于地球表面曲率），这是成立的；大气的“波导”效应在雷达和超地平通信中真的被利用过，只是一般只能多传几百公里，成不了完整的圆。真正拦住它的是：地球大气只有几十公里厚，梯度远不够强；而且大气不均匀、有湍流，光路维持不了那么精致。这个思想实验的价值在于：射线模型里，只要允许 \(n\) 连续变化，“直线”就再也不是特权——光路可以弯成任何形状，全看折射率的分布。这把钥匙留到以后：当引力让时空本身弯曲时，光路的弯曲还会以另一种面目回来（第 41 章）。
\end{thought}

\begin{puzzles}
\item 潜水员仰头看见一个圆形的亮窗，整个岸上的世界——天空、太阳、岸边的树——都被压缩在这个圆窗里，圆窗外是水下景物的镜像。请用临界角说明：为什么岸上四面八方、总共半个球面的景物，能全被塞进一个半顶角只有 \(48.6^\circ\) 的圆锥里？（思路提示：想想岸上贴近地平线方向射来的光，入水时入射角接近 \(90^\circ\)，折射角最大能到多少？斯涅尔定律会把折射角“封顶”在哪里？压缩的不是物体的数量，而是角度。）
\item 钻石的折射率是 \(2.42\)，普通玻璃约 \(1.5\)。为什么切割得当的钻石看起来比玻璃“闪”得多——光一旦进去就迟迟不肯出来？（思路提示：算一算两种材料对空气的临界角。临界角越小，意味着什么角度的入射光会被全反射关在里面？钻石切工的几十个刻面，就是在为光设计一座走不出去的镜子迷宫。）
\item 一根光纤长 \(1\) 千米，光在玻璃芯中以 \(2.0\times10^8\) 米每秒前进。如果光沿轴线直着走，需要多长时间？如果它以“之”字形全反射前进，实际走的路程比轴线长出 \(20\%\)，又需要多长时间？这个时差在多路信号同时到达时会带来什么麻烦？（思路提示：先算直走的时间，数量级是微秒；再乘以 \(1.2\)。两路光到达时间的差，会把高速信号的前后脉冲糊在一起——这就是工程师要限制光纤里光线角度范围的原因。）
\item 费马说光走时间取稳定值的路。可是光在出发时怎么“知道”哪条路最省时间？它难道先把所有路线试一遍吗？请用本章的框架回答：这种说法里，哪些是数学模型的合法语言，哪些是把模型误当成了事实？真正的回答藏在哪一章的图景里？（思路提示：区分“定律的稳定值形式”和“粒子做决策”。第 35 章里，波同时走过所有路径、只有稳定值附近的路径相互加强——那时“比较所有路径”会有一个不拟人化的解释。）
\end{puzzles}

\begin{mainline}
本章我们做了三件事：把光抽象成光线，用两条几何定律（反射、折射）统一了影子、镜面和水面的弯折；用折射率 \(n=c/v\) 把“拐弯”和“变慢”钉成同一件事；又用费马原理把两条定律收成一句“光走时间稳定的路”。下一章，我们给这套几何语言加上镜子和透镜，去回答眼睛、相机和望远镜是怎样成像的；而再往后一章，一束穿过双缝的光将逼着我们把“射线”这张地图换掉——光的另一副面孔，是波。
\end{mainline}
